\documentclass[conference]{IEEEtran}
\IEEEoverridecommandlockouts

% --- Packages ---
\usepackage{cite}
\usepackage{amsmath,amssymb,amsfonts}
\usepackage{algorithmic}
\usepackage{graphicx}
\usepackage{textcomp}
\usepackage{xcolor}
\usepackage{booktabs}
\usepackage{tabularx}
\usepackage{hyperref}
\usepackage{multirow}
\usepackage{float}
\usepackage{algorithm}
\usepackage{enumitem}

\begin{document}

\title{AgriSense AI: A High-Fidelity Hybrid Framework for Multivariate Price Forecasting and Ecological Intelligence}

\author{\IEEEauthorblockN{AgriSense AI Development Team}
\IEEEauthorblockA{\textit{B.Tech Computer Science \& Artificial Intelligence} \\
\textit{Rishihood University}}
}

\maketitle

\begin{abstract}
Agricultural volatility in emerging markets creates significant information asymmetry for smallholder farmers. This research presents AgriSense AI, a unified decision-support platform that integrates high-precision price forecasting with ecological suitability modeling. We introduce a decoupled modular architecture that fuses (i) a Temporal Fusion Transformer (TFT) achieving 96.15\% accuracy on 14-day price windows, (ii) a GOSS-optimized LightGBM point-forecaster, (iii) a high-dimensional soil-climate recommendation engine, and (iv) an interpretability-first plant disease radar. Our results demonstrate that while zero-shot foundation models (Moirai) struggle with localized market shocks, the specialized TFT-LGBM hybrid provides robust quantile-risk corridors for financial stabilization. The system is deployed via a Next.js/Flask dashboard, providing a professional-grade interface for real-time agricultural intelligence.
\end{abstract}

\begin{IEEEkeywords}
Temporal Fusion Transformer, LightGBM, Precision Agriculture, Time-Series Forecasting, Transfer Learning, EfficientNet-B0, AgriSense.
\end{IEEEkeywords}

\section{Introduction}
The Indian agricultural sector, contributing 18\% to GDP, remains one of the most volatile economic environments globally. Commodities like tomatoes and onions often experience 400\% price swings within single harvest cycles \cite{nitiaayog2019}. This research addresses the critical need for forward-looking predictive intelligence by unifying disparate data streams—market prices, NASA weather telemetry, and soil profiles—into a singular, actionable framework.

\section{Methodology Evolution}
Building upon initial explorations in econometrics, AgriSense AI evolved through three distinct phases:
\begin{itemize}[leftmargin=*]
    \item \textbf{Phase 1 (Foundation):} Focused on Agmarknet data scraping and baseline regression models (ARIMA, Random Forest).
    \item \textbf{Phase 2 (Hybridization):} Introduced the LightGBM point-forecaster and initial TFT prototypes, identifying the need for exogenous weather integration.
    \item \textbf{Phase 3 (Production):} Finalized the \textit{Decoupled Engine Pattern}, achieving 96.15\% accuracy on the Epoch 15 TFT model and implementing the unified Next.js dashboard.
\end{itemize}

\section{System Architecture}

\subsection{The Decoupled Engine Pattern}
A core innovation of this project is the modularization of inference. Unlike monolithic scripts, AgriSense AI utilizes a canonical \texttt{Engines/} interface. This allows the backend to dynamically switch between the probabilistic TFT model (for risk analysis) and the deterministic LightGBM model (for market timing).

\subsection{Feature Engineering Density}
We compute over 70 engineered markers to capture market nuances:
\begin{equation}
    F_{seasonal} = \sin\left(\frac{2\pi \cdot d}{365}\right) + \cos\left(\frac{2\pi \cdot d}{365}\right)
\end{equation}
\begin{equation}
    P_{volatility} = \sqrt{\frac{1}{N-1} \sum_{i=1}^{N} (x_i - \bar{x})^2}
\end{equation}
These Fourier transforms and rolling standard deviations allow the models to "see" seasonality and supply shocks before they manifest in price surges.

\section{Model Performance & Evaluation}

\subsection{Price Forecasting Breakthrough}
The transition from Epoch 9 to Epoch 15 of the TFT architecture yielded a massive precision gain. By increasing the depth of the Gated Residual Networks (GRN) and refining the Variable Selection Network (VSN), we achieved:
\begin{itemize}
    \item \textbf{TFT Accuracy (Epoch 15):} 96.15\% (SMAPE: 3.85\%)
    \item \textbf{LightGBM MAE:} 160.41 INR/qtl
\end{itemize}

\subsection{Ecological & Pathological Intelligence}
The **Crop Recommendation** module utilizes 45 composite soil-climate markers, achieving a near-perfect F1-score of 0.99. The **Plant Disease** module, powered by EfficientNet-B0, provides 98.42\% accuracy across 38 classes, with Grad-CAM visualizations identifying biological triggers in real-time.

\section{The Hybrid Advisory Algorithm}
The final advisory output is generated by a multi-step filter that maximizes both biological safety and economic gain:
\begin{enumerate}
    \item \textbf{Biological Filter:} Identify top-3 crops based on soil suitability.
    \item \textbf{Economic Filter:} Run 14-day TFT forecasts for the identified crops in the target mandi.
    \item \textbf{Risk Signal:} Overlay disease radar alerts if the target region shows high humidity/temperature anomalies.
\end{enumerate}

\section{Technocratic Assessment & Failure Analysis}
The project acknowledges the \textit{"Dataset Purity Gap"}. While benchmark accuracies are exceptional, real-world deployment faces sensor noise and "Black Swan" events (unseasonal storms) not captured in historical logs. Future iterations will integrate real-time news sentiment to bridge this gap.

\section{Conclusion}
AgriSense AI provides a robust, production-ready framework for agricultural intelligence. By combining the probabilistic depth of Transformers with the tabular speed of Gradient Boosting, we deliver a system that is both accurate and interpretable, ready for field deployment and collaborative research.

\begin{thebibliography}{99}
\bibitem{nitiaayog2019} NITI Aayog, ``Demand and Supply Projections Towards 2033,'' Government of India, 2019.
\bibitem{lim2021tft} B. Lim et al., ``Temporal Fusion Transformers,'' \textit{IJF}, 2021.
\bibitem{tan2019efficientnet} M. Tan and Q. Le, ``EfficientNet,'' \textit{ICML}, 2019.
\bibitem{ke2017lightgbm} G. Ke et al., ``LightGBM: A Highly Efficient GBDT,'' \textit{NeurIPS}, 2017.
\end{thebibliography}

\end{document}
